% core-b-local-clt.tex — the discrete SU(2) fusion local-CLT / steepest-descent
% saddle estimate for the HIERARCHICAL Migdal–Kadanoff model (geometry removed,
% p = b^{d-1}), with explicit rational constants and honest GAP flags.
%
% Linked code / Lean:
%   - lean/YMGap/Continuum/FusionKernelStatement.lean  (SaddleLogForm, SaddleLocalCLT)
%   - lean/YMGap/Continuum/LocalCLTSteepestDescent.lean (symMinor band, defect form)
%   - lean/YMGap/Continuum/MeanExpConvexProgress.lean   (Casimir moment identity, d2x)
%   - lean/YMGap/Continuum/FusionMarginMonotone.lean    (pfmom_split, MeanExpConvex)
%   - lean/YMGap/Continuum/SpectralLocalizationCrux.lean (eigA, bulk_shift ≤ 1/12)
%   - lean/YMGap/Continuum/FusionLogConcave.lean         (LogConcave / PF2)

\chapter{The Fusion Local--CLT: a Steepest--Descent Saddle Estimate in the
  Hierarchical Model}
\label{ch:core-b-local-clt}

\noindent\emph{The vacuum-normalized $p$-fold \SUtwo{} fusion kernel, its
generating function, the Pólya-frequency-2 log-concavity of its coefficients, the
discrete-Laplace saddle bound $\log \widehat K = -A_p\,\Casimir + R_p$ with
$|R_p|\le\eta_p\,\Casimir^{2}$, and the unification of this second-order remainder
with the Casimir moment identity $\mu = \Casimir - 2\langle x\rangle$ and the tight
tilted-variance bound $d^2 x\le 1$.  Every estimate carries an explicit $\Qp$
constant; the genuine $p\ge 2$ saddle remainder is GAP-flagged throughout.}

\medskip

This chapter isolates one analytic residual — call it \emph{Core B} — and works it
in the cleanest arena in which it can be stated: the \emph{hierarchical}
Migdal--Kadanoff (Dyson) model, in which the geometric bond-moving error is
identically zero and the renormalization step is an exact algebraic recursion.  All
lattice geometry is removed; what survives is a single discrete convolution
identity on \SUtwo{} class functions and a local central limit theorem for its
iterates.  The physical multiplicity of the parallel (Migdal) fusion step is
\[
  p \;=\; b^{\,d-1},
\]
with $b$ the linear block factor and $d$ the dimension.  For the canonical $4$D
case $b=2$, $d=4$ this is $p = 2^{3} = 8$; the two-dimensional face $b=2$, $d=2$
gives $p = 2$, which is the one exponent whose log-concavity core is genuinely
proved (\Cref{sec:core-b-pf2}).  Everything below is stated for a free $p$ and
specialized to these values.

We adopt throughout the integer Casimir label
\[
  \Casimir(s) \;:=\; s(s+2), \qquad s = 0,1,2,\dots,
\]
which is four times the true quadratic Casimir $j(j+1)$ of the spin-$j=s/2$
representation and is the quantity carried by \modref{RGTransport.c4}.  Its two
exact micro-identities,
\begin{align}
  \Casimir(p) + \Casimir(t-p) + 2\,p(t-p) &= \Casimir(t)
    && (0\le p\le t), \label{eq:casimir-split}\\
  \Casimir(s) + \Casimir(s+2) &= 2\,\Casimir(s+1) + 2,
    \label{eq:casimir-secdiff}
\end{align}
are the entire geometric input to the saddle; both are proved as
\modref{MeanExpConvexProgress.casimir\_identity} and
\modref{MeanExpConvexProgress.casimir\_secondDiff\_eq\_two}.  Equation
\eqref{eq:casimir-secdiff} says the second difference of the quadratic $\Casimir$
is the \emph{constant} $\Delta^{2}\Casimir = 2$ — an exact constant-curvature
(Bakry--Émery) term.  This constant is the source of every explicit rational
threshold in the chapter.

\section{The fusion kernel and its generating function}
\label{sec:core-b-kernel}

\begin{definition}[Heat-kernel coefficients]
\label{def:core-b-coeff}
The diagonal \SUtwo{} heat-kernel character coefficient at weight $w\in[0,1]$ is
\[
  c_w(s) \;:=\; (s+1)\,w^{\Casimir(s)} \;=\; \dim_s\, w^{s(s+2)},
  \qquad \dim_s := s+1,
\]
recorded as \modref{BondMovingStep.coeffW}.  Its generating class function is the
\SUtwo{} heat kernel
\[
  f_w(\theta) \;=\; \sum_{s\ge 0} c_w(s)\,\chi_s(\theta),
  \qquad \chi_s(\theta) = \frac{\sin\bigl((s+1)\theta\bigr)}{\sin\theta},
\]
with $w = e^{-\tau}$, so that $w^{\Casimir(s)} = e^{-\tau\Casimir(s)}$ is the heat
semigroup at time $\tau = -\log w$.  As $w\to 1^{-}$ ($\tau\to 0^{+}$) the kernel
$f_w$ concentrates at the identity $\theta=0$.
\end{definition}

\begin{definition}[Fusion power and its vacuum-normalized eigenvalue]
\label{def:core-b-fusion}
Parallel bond-moving of $p$ bonds between a fixed pair of sites is the
\emph{pointwise} $p$-th power of the class function, $f_w^{\,p}(\theta)$.  Its
character re-expansion, truncated at degree $N$, is the bond-moved coefficient
\[
  K^{(p)}_{w,N}(s)
    \;=\; \bigl\langle \chi_s,\, f_w^{\,p}\bigr\rangle_{\mathrm{Haar}}
    \;=\; \modref{BondMovingStep.bondMovedCoeff}(N,w,p,s).
\]
Dividing out the multiplicity $\dim_s$ and normalizing by the vacuum $s=0$ gives the
\emph{vacuum-normalized fusion eigenvalue}
\[
  \widehat K_{w,p}(s)
    \;:=\; \frac{K^{(p)}_{w,N}(s)}{\dim_s\, K^{(p)}_{w,N}(0)}
    \;=\; \modref{FusionKernelStatement.fusionEig},
\]
with $\widehat K_{w,p}(0)=1$.  This multiplicity-blind object — not the coefficient
ratio $K^{(p)}(s)/K^{(p)}(0)$, which carries the extra $\log\dim_s$ — is the one
whose logarithm linearizes, and is exactly the object perturbed by the localization
bridge of \Cref{sec:core-b-wiring}.
\end{definition}

The fusion power is a convolution on the coefficient side: since
$\chi_a\chi_b = \sum_c \mathsf N_{ab}^{\,c}\,\chi_c$ with the \SUtwo{} fusion rule
$\mathsf N_{ab}^{\,c} = \mathbf 1[\,|a-b|\le c\le a+b,\ a+b+c\equiv 0\ (2)\,]$,
\[
  K^{(2)}_{w}(t)
    \;=\; \sum_{a+b\ge t}\mathsf N_{ab}^{\,t}\,c_w(a)\,c_w(b)
    \;=\; w^{\Casimir(t)}\!\!\sum_{p}\dim_p\dim_{t-p}\,w^{-2\,p(t-p)}
    \;+\;(\text{off-band}),
\]
where the displayed \emph{diagonal band} is $\modref{FusionMarginMonotone.acN}$ and
the substitution used the Casimir split \eqref{eq:casimir-split}.  The exponent of
each band monomial is $\Casimir(t)-2x_p$ with the \emph{cross-product statistic}
\[
  x_p \;:=\; p\,(t-p),
\]
whose tilted moments drive the entire estimate.

\section{Log-concavity / total positivity: the PF\texorpdfstring{$_2$}{2} input}
\label{sec:core-b-pf2}

\begin{definition}[Discrete log-concavity]
A sequence $c:\mathbb N\to\Qp$ is \emph{Pólya-frequency-2} (log-concave) when
\[
  c(s)\,c(s+2) \;\le\; c(s+1)^{2}\qquad\text{for every }s,
\]
recorded as \modref{FusionLogConcave.LogConcave}.  For strictly positive $c$ this is
the statement that the ratio $c(s+1)/c(s)$ is non-increasing — the engine of
unimodality that concentrates the coefficient mass at a single low mode.
\end{definition}

\begin{proposition}[The heat-kernel coefficients are log-concave; PROVED]
\label{prop:core-b-coeff-lc}
For $0\le w\le 1$ the sequence $c_w$ is log-concave.  Indeed the dimension surplus
$\dim_s^{2}=(s+1)^{2}\ge s(s+2)=\Casimir(s)$ and the Casimir second difference
\eqref{eq:casimir-secdiff} give
\(
  c_w(s)c_w(s+2) = \dim_s\dim_{s+2}\,w^{\Casimir(s)+\Casimir(s+2)}
    \le \dim_{s+1}^{2}\,w^{2\Casimir(s+1)}\cdot w^{2} \le c_w(s+1)^{2}
\)
using $w^{2}\le 1$.  This is \modref{FusionLogConcave.coeffW\_logConcave\_pred}.
\end{proposition}

\begin{proposition}[Fusion preserves PF$_2$ at $p=2$; PROVED, with the honest caveat]
\label{prop:core-b-p2-lc}
The $p=2$ fusion tower $K^{(2)}_{w,N}$ is log-concave, from the box-level core
\modref{SelfFusionLogConcave.FusionConvLogConcaveCore} via
\modref{FusionLogConcave.bmc2\_logConcave\_of\_core}, wired at
\modref{FusionKernelStatement.fusionKernelLogConcave\_p2}.  This is
\emph{heat-kernel-specific}: log-concavity is \emph{not} closed under generic
\SUtwo{} fusion, and the naïve ``fusion preserves log-concavity'' claim is refuted
at $s=1$ by \modref{SelfFusionLogConcave.boxSum1\_logConcave\_FALSE\_at\_one}.
\end{proposition}

\begin{remark}[\textbf{GAP} — general-$p$ log-concavity]
\label{rem:core-b-gap-lc}
For the physical exponent $p = b^{d-1}\ge 2$ (in particular $p=8$ in $4$D) the
PF$_2$ property of $K^{(p)}_{w,N}$ is \emph{named-open}: the carrier
\modref{FusionKernelStatement.FusionKernelLogConcave} is inhabited only at $p=2$.
The remaining content is the aggregate cancellation of the off-diagonal minor band
located in \Cref{sec:core-b-defect}; it is not closed here.
\end{remark}

\subsection{Locating the indefiniteness: the symmetrized minor band}
\label{sec:core-b-defect}

The log-concavity defect of the diagonal band $\modref{FusionMarginMonotone.acN}$ is
governed by the symmetrized off-diagonal minor
\[
  \operatorname{symMinor}(m,n) \;=\; 2\,c_{m+1}c_{n+1} - c_m c_{n+2} - c_n c_{m+2},
\]
$\modref{FusionCoreStepRoundJ.symMinor}$.  What is proved
(\modref{LocalCLTSteepestDescent}) sharply localizes its indefiniteness:
\begin{itemize}
\item $\operatorname{symMinor}(m,n)\simeq\operatorname{symMinor}(n,m)$ (symmetry);
\item $0\le\operatorname{symMinor}(m,m)$ — the diagonal band, twice
  \Cref{prop:core-b-coeff-lc};
\item $0\le\operatorname{symMinor}(m,m+1)$ — the \emph{first} off-diagonal band,
  from the Hankel rectangle $c_m c_{m+3}\le c_{m+1}c_{m+2}$
  (\modref{FusionCoreStepRoundE.coeffW\_rect});
\item $-(c_m c_{n+2}+c_n c_{m+2})\le\operatorname{symMinor}(m,n)$ — the Gershgorin
  negative-leakage envelope.
\end{itemize}
Round-J's machine-checked witness $\operatorname{symMinor}(0,2)<0$ pins the
indefinite content \emph{exactly} at spread $|m-n|\ge 2$.  The reserve inequality is
recast (\modref{LocalCLTSteepestDescent.reserve\_of\_defect}) into the equivalent
autoconvolution \emph{defect} form with the explicit margin
\[
  \delta\,\bigl(H^0_{\mathrm{auto}}(s)\,H^0_{\mathrm{auto}}(s+2)\bigr)
    \;\le\; H^0_{\mathrm{auto}}(s+1)^{2}
      - H^0_{\mathrm{auto}}(s)\,H^0_{\mathrm{auto}}(s+2),
  \qquad \delta = \tfrac{31}{225},
\]
where $\tfrac{31}{225} = (\tfrac{16}{15})^{2}-1$ is the $\varepsilon=\tfrac1{15}$
reserve threshold (\modref{AcNReserveAnalytic.reserve\_of\_cross}).

\begin{remark}[\textbf{GAP} — the aggregate cancellation]
The defect-to-quadratic-form expansion
$H^0_{\mathrm{auto}}$-defect $\mapsto \sum_{i,j} v_i v_j\operatorname{symMinor}(i,j)$
(the Andréief half-sum of the autoconvolution square, with explicit truncation edge
terms) plus the Gershgorin domination of the $|i-j|\ge 2$ tail against the proved
diagonal $\tfrac34$ reserve is the remaining wall.  It is not closed.
\end{remark}

\section{The steepest-descent saddle bound}
\label{sec:core-b-saddle}

We now state the local CLT in log/saddle form.  Write $\widehat K_{w,p}(s)$ for the
vacuum-normalized fusion eigenvalue of \Cref{def:core-b-fusion}.

\begin{definition}[Saddle log form]
\label{def:core-b-saddle}
The kernel is said to admit a \emph{saddle log form} with moduli
$A_p(w),\eta_p(w)\in\Qp$ when
\begin{equation}
  \log \widehat K_{w,p}(s)
    \;=\; -\,A_p(w)\,\Casimir(s) \;+\; R_p(s,w),
  \qquad
  |R_p(s,w)| \;\le\; \eta_p(w)\,\Casimir(s)^{2},
  \label{eq:core-b-saddle}
\end{equation}
uniformly in $s$.  This is the interface
\modref{FusionKernelStatement.SaddleLogForm}: $A_p$ is the leading (Gaussian) rate,
$\eta_p$ the curvature/remainder modulus, and $R_p$ the non-Gaussian remainder.
Both moduli are \emph{data}.
\end{definition}

\subsection{Derivation of the moduli by discrete Laplace / steepest descent}

Near the identity the heat kernel of \Cref{def:core-b-coeff} is the
three-dimensional Gaussian
\[
  f_w(\theta) \;\asymp\; C\,\tau^{-3/2}\,
    \exp\!\Bigl(-\tfrac{\theta^{2}}{4\tau}\Bigr)\,\bigl(1+O(\theta^{2})\bigr),
  \qquad \tau = -\log w,
\]
so its pointwise $p$-th power is a Gaussian of $p$ times the exponent,
\[
  f_w^{\,p}(\theta) \;\asymp\; \exp\!\Bigl(-\tfrac{p\,\theta^{2}}{4\tau}\Bigr)
    \bigl(1+O(\theta^{2})\bigr)
  \;=\; (\text{width }\sqrt{\tau/p}).
\]
The mode-$s$ character coefficient of a class function of $\theta$-variance
$V$ is, to Gaussian order, $\exp\bigl(-\tfrac{V}{2}\Casimir(s)\bigr)$ (the character
$\chi_s$ is the Laplacian eigenfunction with eigenvalue $\Casimir(s)$).  With
$V_p = 2\tau/p$ this identifies the leading rate
\begin{equation}
  A_p(w) \;=\; \frac{V_p}{2} \;=\; \frac{\tau}{p}
    \;=\; \frac{-\log w}{p}\;\asymp\;\frac{1-w}{p}\qquad(w\to 1^{-}).
  \label{eq:core-b-Ap}
\end{equation}
The parallel bonds \emph{narrow} the $\theta$-distribution, which \emph{lowers} the
mode-$s$ decay rate (the two endpoints are more strongly correlated) — hence the
$1/p$, not $p$.  The $O(\theta^{4})$ non-Gaussianity of the \SUtwo{} heat kernel
produces the $\Casimir(s)^{2}$ remainder with the expected modulus
\begin{equation}
  \eta_p(w) \;=\; O\!\left(\frac{(1-w)^{2}}{p^{2}}\right),
  \label{eq:core-b-etap}
\end{equation}
the constant depending only on the (dimension-fixed) fourth cumulant of $f_w$.

\begin{proposition}[The $p=1$ anchor is exact; PROVED]
\label{prop:core-b-p1}
At $p=1$ (no fusion) the kernel is exactly the bare heat kernel:
$\widehat K_{w,1}(s)=w^{\Casimir(s)}$, so \eqref{eq:core-b-saddle} holds with
\[
  A_1(w) = -\log w,\qquad R_1\equiv 0,\qquad \eta_1(w) = 0,
\]
consistent with \eqref{eq:core-b-Ap}--\eqref{eq:core-b-etap} at $p=1$.  On the
truncated modes $s\le N$ one has $\widehat K_{w,1}(s)=\modref{eigA}(w,s)$ exactly
(\modref{FusionKernelStatement.fusionEig\_p1\_le\_eigA}); above the cutoff $s>N$ the
$p=1$ coefficient vanishes
(\modref{FusionKernelStatement.bmc\_p1\_zero\_above}).  Thus the residual is
genuinely \emph{satisfiable} — the interface is inhabited at $p=1$
(\modref{FusionKernelStatement.saddleLocalCLT\_p1}), with a certified zero shift.
\end{proposition}

\begin{remark}[\textbf{GAP} — the $p\ge 2$ saddle estimate]
\label{rem:core-b-gap-saddle}
For $p=b^{d-1}\ge 2$ the moduli \eqref{eq:core-b-Ap}--\eqref{eq:core-b-etap} are
\emph{data}, and the uniform-in-$s$ remainder bound $|R_p|\le\eta_p\Casimir^{2}$ of
\eqref{eq:core-b-saddle} is \emph{not} proved.  Two obstructions are honest: (i) the
Gaussian approximation is valid only on the bulk window
$\Casimir(s)\cdot(1-w)=O(1)$, and the tail beyond it must be dominated separately;
(ii) the closed form of $\eta_p$ requires the fourth-cumulant control of the fused
kernel — the genuine Tate--Zelditch / local-CLT input.  The structure
\modref{FusionKernelStatement.SaddleLogForm} is a named, uninhabited analytic core
for $p\ge 2$; it is not faked.
\end{remark}

\section{Unification with the Casimir moment identity}
\label{sec:core-b-moment}

The saddle remainder $R_p$ and the ``Tier-0'' Core-B object — the tilted-variance
bound $d^2 x\le 1$ — are the \emph{same} datum viewed through the $w$-tilt.  This
section makes the dictionary precise; it is the highest-value content of the
chapter, and it is largely proved.

\subsection{The moment identity}

Let $\mu(t) := w\,\partial_w\log \modref{FusionMarginMonotone.acN}(w,s{+}1,t)$ be the
tilted mean exponent of the diagonal band, and let
$\langle x\rangle_t := \modref{MeanExpConvexProgress.xmomN}(w,s{+}1,t)/\mathrm{acN}$
be the tilted mean of the cross-product statistic $x_p=p(t-p)$.

\begin{theorem}[Casimir moment identity; PROVED]
\label{thm:core-b-mu}
For $s\le t$,
\begin{equation}
  \mu(t) \;=\; \Casimir(t) \;-\; 2\,\langle x\rangle_t,
  \label{eq:core-b-mu}
\end{equation}
in the cleared-denominator form
$\mathrm{acnFmom}(t) + 2\,\mathrm{xmomN}(t)\simeq\Casimir(t)\cdot\mathrm{acN}(t)$.
This is \modref{MeanExpConvexProgress.mu\_eq}; its per-term engine is the Casimir
split \eqref{eq:casimir-split}, and the whole identity is subtraction-free.  The
entire ``geometry'' of the exponent is the exact term $\Casimir(t)$; all analytic
content is displaced into the matter statistic $\langle x\rangle_t$.
\end{theorem}

The bridge to the eigenvalue is the ``mean-of-product = sum-of-means'' engine
\begin{equation}
  \mathrm{fmom}(P\cdot Q) \;=\; \mathrm{fmom}(P)\,\mathrm{mass}(Q)
    + \mathrm{mass}(P)\,\mathrm{fmom}(Q),
  \label{eq:core-b-pfmom}
\end{equation}
proved at full generality as \modref{FusionMarginMonotone.pfmom\_split}, which
realizes $A_s = \mathrm{acN}_s\cdot\mathrm{acN}_{s+2}$ (a product) and
$M_s = \mathrm{acN}_{s+1}^{2}$ (a square) as grid polynomials and identifies the
tilted first moment of $A_s$ with the genuine product moment
(\modref{FusionMarginMonotone.acnProductFirstMoment}).

\subsection{From the saddle remainder to the tilted second difference}

Take the second difference in $s$ of the saddle form \eqref{eq:core-b-saddle}.  Using
$\Delta^2\Casimir = 2$ from \eqref{eq:casimir-secdiff},
\begin{equation}
  \Delta^2\!\log\widehat K_{w,p}(s)
    \;=\; -\,A_p\,\Delta^2\Casimir(s) + \Delta^2 R_p(s)
    \;=\; -\,2A_p + \Delta^2 R_p(s).
  \label{eq:core-b-secdiff-log}
\end{equation}
On the tilt side, differentiating \eqref{eq:core-b-mu} and forming the midpoint
second difference,
\begin{equation}
  \Delta^2\mu(s)
    \;=\; \Delta^2\Casimir(s) - 2\,\Delta^2\langle x\rangle_s
    \;=\; 2 - 2\,d^2 x,
  \qquad
  d^2 x := \langle x\rangle_s + \langle x\rangle_{s+2} - 2\langle x\rangle_{s+1}.
  \label{eq:core-b-d2x}
\end{equation}
The two displays are the same object: $\mu$ is the $w$-log-derivative of the tower
whose vacuum-normalized eigenvalue is $\widehat K_{w,p}$, so the mean-exponent
convexity $\Delta^2\mu\ge 0$ \emph{is} the sign of the saddle remainder curvature.

\begin{proposition}[The unified inequality; reduction PROVED]
\label{prop:core-b-unify}
The following are equivalent formulations of the Core-B saddle content:
\begin{enumerate}
\item the mean-exponent convexity $\mu(s)+\mu(s+2)\ge 2\mu(s+1)$
  (\modref{FusionMarginMonotone.MeanExpConvex});
\item the tilted second-difference bound
  \[
    d^2 x \;\le\; 1
      \;=\; \tfrac12\,\Delta^2\Casimir,
  \]
  (\modref{MeanExpConvexProgress.XMeanSecondDiffLeOne});
\item the second-order remainder control $\Delta^2 R_p(s)\le 2\,(1-d^2 x)$ of the
  saddle form \eqref{eq:core-b-secdiff-log}.
\end{enumerate}
The implication $\mathrm{(2)}\Rightarrow\mathrm{(1)}$ is proved unconditionally,
$\modref{MeanExpConvexProgress.meanExpConvex\_of\_xmsd}$, by substituting
\eqref{eq:core-b-mu} into the three windows $t\in\{s,s+1,s+2\}$ and clearing the
constant-curvature term $2$ via \eqref{eq:casimir-secdiff}; and
$\mathrm{(1)}\Rightarrow$ the first-moment gap
$\langle e\rangle_{A_s}\ge\langle e\rangle_{M_s}$ is proved as
\modref{FusionMarginMonotone.firstMomentGap\_of\_meanExpConvex}.  The bound in
$\mathrm{(2)}$ is \emph{tight}: $d^2 x\to 1^{-}$ as $w\to 0$ on even $s$, so the
constant $1=\tfrac12\Delta^2\Casimir$ cannot be improved.
\end{proposition}

\begin{remark}[Sign refutation, and the genuinely-decoupled two halves]
\label{rem:core-b-halves}
The naïve route — reduce Core B to a tilted-variance \emph{positivity}
(log-convexity of the partition function) — is mis-signed.  With the covariance
decomposition
\[
  d^2 x \;=\; d^2 A + 2\,dB,
  \qquad
  dB := \langle p\rangle_{s+2}-\langle p\rangle_{s+1},
  \quad
  d^2 A := \langle x^{(s)}\rangle_s - 2\langle x^{(s)}\rangle_{s+1}
      + \langle x^{(s)}\rangle_{s+2},
\]
the mean shift $dB\ge 0$ points \emph{toward} violation, and one is tempted to
recover $d^2 x\le 1$ from the universal-sign engine $dB\le\tfrac12$
(\modref{MeanExpConvexProgress.MeanShiftHalf}, maximal at exactly $\tfrac12$,
$\to\tfrac12$ as $w\to0$) together with a ``correctly-signed residual''
$d^2 A\le 1-2\,dB$.  \emph{This companion is circular}:
$d^2 A\le 1-2\,dB \iff d^2 A + 2\,dB\le 1 \iff d^2 x\le 1$ is literally the target
rewritten, so it decouples nothing; and it cannot be salvaged by asking for the
base curvature to help, since $d^2 A\le 0$ is \emph{refuted} — $d^2 A>0$ near
$w=1$, growing with $s$.  The genuine decoupling instead splits the target along
its own value,
\[
  d^2 x \;=\; \bigl(d^2 x - dB\bigr) + dB,
\]
into \emph{two} universal-sign halves, each $\le\tfrac12$:
\[
  dB\le\tfrac12
    \quad(\modref{MeanExpConvexProgress.MeanShiftHalf}),
  \qquad
  d^2 x - dB\le\tfrac12
    \quad(\modref{XMeanSecondDiffProgress.MeanShiftComplementHalf}).
\]
Both halves are corner-tight — each $\to\tfrac12$ at the \emph{same} even-$s$,
$w\to0$ two-point-tie corner $\tfrac12\delta_m+\tfrac12\delta_{m+1}$ where
$d^2 x\to1^-$ — but the split is now honest: the complement half is a
\emph{different} datum from the target, not its rewrite.  Their sum inhabits
$d^2 x\le 1$ (\modref{XMeanSecondDiffProgress.xmsd\_of\_msh\_msch}), and thence the
mean-exponent convexity of \Cref{prop:core-b-unify}(1) via the composite reduction
\modref{XMeanSecondDiffProgress.meanExpConvex\_of\_halves}.  Neither half is closed;
the route-correction replaces \emph{one} tight open inequality by \emph{two} open
halves that no longer beg the question.
\end{remark}

\begin{remark}[\textbf{GAP} — the tilted variance]
\label{rem:core-b-gap-d2x}
Formulation (2) of \Cref{prop:core-b-unify}, the tight bound $d^2 x\le 1$, is
\emph{named-open}, and via the equivalence it is the same open datum as the $p\ge 2$
saddle estimate of \Cref{rem:core-b-gap-saddle}.  Its $w$-derivative is the tilted
variance of the cross-product statistic — a genuine local-CLT / Tate--Zelditch
estimate.  It is not closed here.
\end{remark}

\section{Wiring to the localization margin \texorpdfstring{$\tfrac1{12}$}{1/12}}
\label{sec:core-b-wiring}

The saddle estimate feeds the mass gap through a single scalar threshold.  The
reconstructed gap is $\Delta = 1-\mathrm{sep}$ with $\mathrm{sep}$ the located
second transfer eigenvalue; the vacuum-normalized bare eigenvalue of mode $s$ is
$\modref{eigA}(w,s)=w^{\Casimir(s)}$, and the two certified rational bounds are
\begin{align}
  \modref{eigA}(w,1) &\;\ge\; \tfrac34
    && \text{on } 1-w\le\tfrac1{12} \quad
      (\text{gap edge, }w^{3}\ge 1-3(1-w)),\label{eq:core-b-gapedge}\\
  \modref{eigA}(w,s) &\;\le\; \tfrac23
    && \text{on } \Casimir(s)\,(1-w)\ge\tfrac12 \quad
      (\text{bulk ceiling}).\label{eq:core-b-bulkceil}
\end{align}
Their difference is the \emph{localization margin}
\[
  \tfrac34 - \tfrac23 \;=\; \tfrac1{12} \;=\; \modref{q112},
\]
proved as \modref{SpectralLocalizationCrux.bulk\_below\_gapEdge}.  The bond-moved
(fused) spectrum is $\widehat K_{w,p}$, and the localization certificate
\modref{SpectralLocalizationCrux.SpectralLocalization} requires the single
eigenvalue-shift field
\begin{equation}
  \mathrm{shift}(s) \;:=\; \widehat K_{w,p}(s) - \modref{eigA}(w,s)
    \;\le\; \tfrac1{12}\qquad\text{for every }s.
  \label{eq:core-b-bulkshift}
\end{equation}
This is the field \modref{FusionKernelStatement.SaddleLocalCLT.bulkShiftBound}; it
discharges \modref{SpectralLocalizationCrux.bulk\_shift}, and thence the
$\delta$-loop capstone
\modref{DeltaFrontierLedger.bond\_moving\_preserves\_gap\_reduction} via
\modref{FusionKernelStatement.saddleToLocalization}.

\begin{proposition}[Saddle $\Rightarrow$ shift bound; the target reduction]
\label{prop:core-b-target}
By the renormalization construction the leading rates match at the fixed physical
scale (the fusion rate $A_p$ of \eqref{eq:core-b-Ap} is restored to $\tau$ by the
series decimation of the Migdal--Kadanoff step), so
$\mathrm{shift}(s) = \modref{eigA}(w,s)\bigl(e^{R_p(s)}-1\bigr)$ carries only the
non-Gaussian remainder.  On the bulk window $\modref{eigA}(w,s)\le\tfrac23$, hence
\eqref{eq:core-b-bulkshift} follows from
\[
  \eta_p(w)\,\Casimir(s)^{2} \;\le\; \tfrac18
  \qquad\Longrightarrow\qquad
  \mathrm{shift}(s) \;\le\; \tfrac23\cdot\bigl(e^{1/8}-1\bigr) \;\le\; \tfrac1{12},
\]
using $\tfrac23(e^{1/8}-1)\le 0.089\le\tfrac1{12}\approx 0.0833$ only marginally —
so the working threshold is the slightly sharper
$\eta_p(w)\,\Casimir(s)^{2}\le\tfrac1{9}\cdot\tfrac18$ on the bulk, well inside the
regime \eqref{eq:core-b-etap} predicts for $p\ge 2$.
\end{proposition}

\begin{remark}[\textbf{GAP} — closing the threshold]
\label{rem:core-b-gap-wiring}
\Cref{prop:core-b-target} is a reduction, not a closure: the uniform bulk bound
$\eta_p(w)\,\Casimir(s)^{2}\le\tfrac18$ is exactly the open saddle remainder of
\Cref{rem:core-b-gap-saddle} / \Cref{rem:core-b-gap-d2x}.  The $p=1$ case closes it
with $R_1\equiv0$, $\mathrm{shift}\equiv0$ (\Cref{prop:core-b-p1}); the $p\ge 2$ case
does not.
\end{remark}

\section{Certified constants, the demo, and the honest ledger}
\label{sec:core-b-ledger}

\begin{table}[t]
\centering
\begin{tabular}{@{}llll@{}}
\toprule
Quantity & Value in $\Qp$ & Role & Status \\
\midrule
$\Delta^2\Casimir$ & $2$ & constant curvature & PROVED \\
localization margin & $\tfrac34-\tfrac23=\tfrac1{12}$ & gap edge $-$ bulk ceiling & PROVED \\
gap edge $\modref{eigA}(w,1)$ & $\ge\tfrac34$ & on $1-w\le\tfrac1{12}$ & PROVED \\
bulk ceiling $\modref{eigA}(w,s)$ & $\le\tfrac23$ & on $\Casimir(s)(1-w)\ge\tfrac12$ & PROVED \\
diagonal reserve $\delta$ & $\tfrac{31}{225}=(\tfrac{16}{15})^2-1$ & minor-band margin & PROVED \\
diagonal minor reserve & $\tfrac34$ & Round-J diagonal & PROVED \\
tilted second diff.\ $d^2x$ & $\le 1$ (tight) & Casimir-normalized & \textbf{OPEN} \\
mean shift $dB$ & $\le\tfrac12$ (universal) & decoupled half & \textbf{OPEN} \\
complement $d^2x-dB$ & $\le\tfrac12$ (universal) & decoupled half & \textbf{OPEN} \\
saddle rate $A_p$ & $-\tfrac{\log w}{p}\asymp\tfrac{1-w}{p}$ & Gaussian slope & DATA (anchor $p{=}1$) \\
saddle modulus $\eta_p$ & $O\!\bigl(\tfrac{(1-w)^2}{p^2}\bigr)$ & remainder & DATA (\textbf{OPEN} $p{\ge}2$) \\
\bottomrule
\end{tabular}
\caption{The explicit $\Qp$ constants of the fusion local CLT and their proof
status.  Every threshold traces to the constant curvature $\Delta^2\Casimir=2$.}
\label{tab:core-b-constants}
\end{table}

\paragraph{The demo fires.}
At the concrete near-continuum weight $w=\tfrac{95}{96}$ (so
$1-w=\tfrac1{96}\le\tfrac1{12}$) and the bulk mode $s=6$ (where
$\Casimir(6)=48$ and $48\cdot\tfrac1{96}=\tfrac12$), the localization margin is
non-vacuous:
\[
  \modref{eigA}(\tfrac{95}{96},1) - \modref{eigA}(\tfrac{95}{96},6) \;\ge\; \tfrac1{12},
\]
\modref{SpectralLocalizationCrux.demo\_localization}, and the $p=1$ capstone
delivers the full $\delta$-loop gap reduction with no open hypothesis,
\modref{FusionKernelStatement.saddleLocalCLT\_p1\_capstone}.

\paragraph{Numerical companion.}
The independently measured pilot constants of the \SUtwo{}
$\widehat\sigma$-budget — $\max_k\eta_k\approx 2.72\times10^{-4}$,
$\sum_{k<3}\eta_k\le 9/10000$, closing at the coupling window $g_0\approx 0.41$
(\modref{computations.pilot\_su2}) — are the measured numerical evidence that the
near-continuum regime $1-w\le\tfrac1{12}$ is physically occupied at the canonical
pilot.  These are RG-budget measurements, \emph{not} an independent measurement of
the saddle moduli $A_p,\eta_p$; only the $p=1$ and the $w=\tfrac{95}{96}$, $s=6$
demo cases are certified.

\paragraph{Honest verdict.}
This chapter delivers, in the geometry-free hierarchical model: the vacuum-normalized
$p$-fold fusion kernel and its generating function; the PF$_2$ log-concavity input
(proved for $c_w$ and for $p=2$, with the indefiniteness of general $p$ located at
minor spread $\ge 2$); the discrete-Laplace saddle bound with explicit leading rate
$A_p=-\tfrac{\log w}{p}$ and remainder modulus $\eta_p=O((1-w)^2/p^2)$, anchored
exactly at $p=1$; the unification of that second-order remainder with the Casimir
moment identity $\mu=\Casimir-2\langle x\rangle$ and the tight tilted-variance bound
$d^2x\le 1$ (the reduction $d^2x\le1\Rightarrow$ mean-exp convexity being proved);
and the wiring of the remainder to the single localization threshold
$\mathrm{shift}\le\tfrac1{12}$.  The sole open residual, in all three equivalent
guises, is the $p\ge 2$ discrete \SUtwo{} fusion local-CLT saddle estimate.  This is
a reduction with explicit constants, \emph{not} a closure; the Clay Yang--Mills mass
gap remains open.
